\section{Discussion}
\label{chap:discussion}

In Part \ref{part:results} the results from testing Hover on a selection of users were presented. This chapter will compare and discuss those findings, revealing elements of the study necessary for answering the research questions. The described knowledge will help us conclude the project in the next chapter.

\subsection{Discussing Results on Player Engagement}
Chapter \ref{chap:player_engagement} presented results on player engagement collected from the questionnaire and application usage data. The screen engagement data revealed that the feed, tracking screen, and leaderboard are where the participants spent the most time. This result was expected as these screens hold the core game design elements and interaction mechanisms. The screen engagement measurement only considers the time where the application is foregrounded. Since the gameplay in Hover is passive (tracking typically runs in the background), it is not surprising that the engagement in the tracking screen has approximately the same fraction of use as the feed and leaderboard. 

Though screen engagement reported that the participants spent most of the time in the feed, we found that the participants did not use the reaction functionality a lot. 58\% of participants disagreed that they often reacted to other players' activities and achievements (S5). One reason for this is that the users refrained from interacting with participants who did not have a screen name they recognized. Because of privacy concerns raised by NSD, we could not encourage the tester to use their real names. Consequently, the participants were free to choose their screen names, affecting their willingness to interact with other players. 

Moreover, the prototype did not include the concept of friends or groups. All players could, therefore, potentially interact with everyone. This situation likely influenced the number of interactions because the players felt uncomfortable interacting with people they did not know.

The student society group reacted on average three times as much as the group of classmates. We think the primary reason for the difference is that the student society group has a closer bond. The relationship between the participants in the group affects their willingness to interact with each other in the game.

Out of the 141 activities started, 19 were discarded, and six were lost. The number of discarded activities was not surprising compared to the number of participants in the experiment. Several scenarios could result in a participant not wanting to publish their activity. One example is that the participant started an activity only out of curiosity when exploring the application. The six activities lost are caused by user errors or technical issues with the device or application. However, the number of lost activities is not remarkable considering it only makes up 4\% of the total activities. The 116 published activities give an average of 6 activities per participant, which is roughly equal to completing an activity every other day. If we consider the circumstances, this is a solid number.

Usage data also revealed that the participants played more during the first week (76 activities) of testing than the second week (40 activities). This difference is likely because players were more eager to play initially, and then their excitement decreased in the second week. The decrease was somewhat unexpected, and we had hoped the numbers would increase in the second week instead. The numbers also show a reduction in the activity count during the weekends. This pattern was also expected due to the national guidelines encouraging people to stay at home.

Out of the 46 possible Hover locations, only five were used. This result was somewhat disappointing but expected due to the national regulations and recommendations at the time. Four of the five locations are popular spots for students in Trondheim, and therefore the small diversity among locations was not unexpected considering the test population. None of the participants visited any of the social category locations, e.g., restaurants and cafes. However, this was not surprising as many of these were either closed or restricted to some degree.

The application also gave participants location-based notifications reminding them to hover when they were inside a zone. These notifications positively affected the engagement in the game, with 47\% (S6) agreeing that notifications made them use the application more. Several of the participants stated that they appreciated the reminders. We consider the number of participants agreeing to S6 to be a high share (46\%), but we expected it to be higher. Some participants did not receive location reminders due to incorrect location permissions, despite being repeatedly instructed to provide them. The lack of permission likely affected the engagement, as some participants stated that it would have made them play the game more.

The in-game challenges had a low engagement, with only four challenges created by the participants. For that reason, we found it necessary to create three challenges to make sure all participants were familiar with the feature. There was a significant difference between the two groups, with 86\% of the classmates and only 42\% of the student society group agreeing to the statement. However, we did not see a distinct difference in actual challenge engagement. Thus, it is probable that the classmates’ responses are affected by the familiarity bias. Furthermore, the impression of challenge engagement compared with the responses from the statements is somewhat conflicting. It is not clear why the participants did not engage in this feature more. However, this relates to motivation.

\textbf{Difference between the two groups} \\
There was a remarkable difference in the average score accumulated by the two groups (Avg S = 5640 points, Avg C = 1563 points). However, the average time spent hovering did not differ much (Avg S = 10h:52m points, Avg C = 10h:14m). Thus, the difference in the average score is influenced by something other than the time spent hovering. One apparent reason is that the student society tracked 65\% of their activities in the cultural category (5 points/minute), whereas the classmates tracked 91\% of their activities in the educational category (1 point/minute). Another reason is that the participants from the student society used the friend tracking feature more, resulting in frequent double points. 

There is also a considerable difference in the number of achievements earned by the two groups. The student society group earned almost three times as many achievements as the group of classmates. This difference is reasonable because 15 of the 19 rewards are directly related to the player’s points. Besides, the student society accumulated on average 3.6 times more points than the group of classmates. This difference does not seem to affect the participants’ engagement, with 47\% disagreeing that they were curious about the achievements they could achieve (S1).

\subsection{Discussing Results on Player Motivation}
Chapter \ref{chap:player_motivation} presented the results related to the players’ motivation from the testing of Hover. The questionnaire revealed that the leaderboard, double points, and other social elements increased the players’ motivation to use the application. The social aspect of Hover was a large motivation factor for a lot of the players. This finding was expected, as social elements are common in a lot of applications today. 

Additionally, the interviews also revealed points as a central mechanism for motivating players. As described in the prestudy (Chapter \ref{chap:reward_systems}), points reward systems are common game design elements. Thus, points motivating the players was an expected result.

On the other hand, the results also described how some game design elements did not increase the motivation for the players. The results from the questionnaire detailed that \textit{achievements}, \textit{notifications}, \textit{streaks}, and \textit{reactions} did not impact the players’ motivation for most of the participants.

The low score of achievements was a bit unexpected. Similar to points, achievements are a common element in game reward systems. However, according to the results from the questionnaire, this element scored low on increasing the motivation for the players. It is hard to think of a specific reason for this answer, but it is probably a combination of many factors. Maybe there were too few achievements, or perhaps they were too obvious. Further research and expansion of the achievement system should be done to improve this part of the game.

Notifications also scored unexpectedly low on motivation, despite notifications being a common element in many applications. However, because notifications are such a natural part of other applications, people are used to getting lots of them daily. Consequently, a few notifications from Hover may not motivate the players as much as we expected.

Some of the other feedback was more expected, based on our observation during the testing period. The small number of participants impacted by streaks could have something to do with engagement during the test period. Many players did not post activities regularly and therefore did not experience any streaks at all. A similar implementation of this game design element in Snapchat is proven to increase demand, thus motivation in the social networks \citep{hristova_snapchat_2019}.

The last game design element with disappointing results was reactions. As explained in the engagement discussion above, Hover lacks a proper social network. This shortcoming is likely the reason for the negative feedback to this game element.

The results in Chapter \ref{chap:player_motivation} also revealed that two of the questions had a conflicting result. In question Q17, 84\% of the participants answered that they were never or only a few times motivated by challenges to use Hover. While in S14: \textit{Challenges between friends inspired me to use the application}, 74\% agreed that challenges between friends motivated them. It is impossible that challenges motivate a player if they have never used them. In retrospect, we find the formulation of S14 is slightly wage and we think that some of the participants might have misunderstood the statement referring to the in-game challenges. The answers from Q17 also agree more with other numbers related to the in-game challenges presented in Chapter \ref{chap:player_engagement}.

In the last question (Q19), one-fourth of the participants thought the points correct in relation to their effort in the game. The problem was mainly related to getting more points in the culture than in the educational category. The culture category was meant for theater, musicals, and concerts; thus, the point system was designed with this in mind. We somehow missed how the volunteers at such places naturally would spend a lot of time there. The algorithm for calculating points could have been changed during the test period. But we decided that it was fairest to let it stay as is.

\textbf{Difference between the two groups} \\
As presented in the results, two of the statements in the questionnaire had a statistically significant difference in the answers between the group of classmates and the student society group:

In S11: \textit{“Reactions to my activities inspired me to use the application more.”}, the people in the student society group disagreed while the classmates agreed. We think that this result is related to the familiarity bias since the participants in the student society group used the feature more than the classmates (presented in Chapter \ref{chap:player_engagement}). The classmates may have been a little bit too kind with their answers to this question. 

In S13: \textit{“The ability to earn double points by using the application with a friend inspired me to use it more.”}, the student society group was more positive than the group of classmates. We think this difference is based on the group composition and how the two groups meet. As explained earlier (Chapter \ref{chap:test_population}), the student society group meets regularly. Consequently, their possibility to hover together increases and becomes more natural. Many of the participants in the group of classmates did not meet during the two weeks of testing.

\subsection{Discussing Results on Player Enjoyment}
Results related to player enjoyment were presented in Chapter \ref{chap:player_enjoyment}. The findings were generally good, indicating that many of the participants enjoyed the game. We were pleased to see that the majority of participants confirmed that they liked the concept (S17). Certain game elements had a clear positive influence on the player’s enjoyment. One example is the Hover with Friends functionality, which participants described as fun in the interviews.

However, as expected, some of the answers did not reflect the general response. In S18, users disagreed with \textit{“It was sad to stop using Hover after the test period ended.”} We think the reason is related to the formulation of the statement and that the result could have been different if the question was asked in another way. Something similar to “I missed using Hover after the test period ended” would have been better.

In S23, most participants disagreed that they “felt attached to their avatar.” This result was expected since the functionality was not implemented. FR28 in Chapter \ref{chap:requirements} specified that: \textit{“Users should be able to evolve a character based on their score.”} Because this functional requirement was low-priority, the feature was dropped from our implementation plan. A randomized avatar picture was used instead (illustrated in Figure \ref{fig:avatar_collage} in Chapter \ref{chap:game_design}). Despite this fact, we still wanted to include this question in the survey to get some feedback on how such an implementation could work. However, in retrospect, we regret that we did not alter the question slightly to get more direct feedback. 

The last statement where the answers diverged from the general direction was in S25: \textit{“I was curious to see what my friends had been up to.”} The result from this statement was split between the participants. Roughly half of the respondents disagreed, while the other half agreed. The answer also showed a statistically significant difference between the two groups. It is challenging to explain this difference because it could be complex or as simple as the classmates being more curious than the student society members. Based on this single question, we cannot conclude whether or not the participants’ curiosity improved their enjoyment.

As the last question in this part, the participants were asked if they had any technical problems during the test period. A bit surprisingly, more than  half of the respondents answered “\textit{yes}.” However, the problem was only related to the Android implementation of the game. Even though we did not expect this result, we recognize how it could have been avoided. During the development of Hover, our approach involved primarily testing the application on Android and iOS emulators and physical iOS devices. Since both of us use an iPhone, it was natural to use these as test devices. At this time, the test participants were not confirmed, and thus we did not know that the majority would use the application on Android. Based on our testing, we knew that there were some differences between the two operating systems, but we did not anticipate that the issues would impact the user experience. In retrospect, we realize that the application should have been tested more on physical Android devices.

\textbf{Difference between the two groups} \\
Answers in four of the statements revealed statistically significant differences between the two groups. The responses to S21, which asked the participants if their status in the game increased the more they played, had a very low p-value of 0.0064. We think that the reason for this difference is that the student society published more activities per person than the classmates did (see results in Chapter \ref{chap:player_engagement}). Since the number of points does not increase with experience, it might feel like your status does not increase as much as it did initially.

S18 and S23 also show a significant difference. There is no apparent reason for the difference, but we think it might be because of familiarity bias. As explained above, we are closer with our classmates than the student society group. Thus, it may be easier for the participants from the student society group to be more transparent and honest in their answers.

\subsection{Discussing Results on Change in Habits}
The last research question, RQ5, investigates how the players’ habits were affected by their use of Hover in the test period. The answers in the questionnaire were generally negative, with most participants disagreeing on all questions. Even though we would have liked to see another result, it is not unexpected. Forming and changing people’s habits takes a long time. Usually, considerably longer than two weeks \citep{lally_how_2010}. Thus, to test if someone has changed their habits appropriately, the testing period should last more than two weeks. Unfortunately, we did not have time for an extended test period in this project.

Q22 asked the participants: \textit{“Would you like to continue using a game like this?”} The answers were split, with around half of the respondents answering “\textit{yes}” and the other half “\textit{no}.” The group of classmates was generally more positive. A bit unexpected since the participants from the student society used the application most. The answers were somewhat disappointing, and we had hoped to see better results. However, it was likely influenced by the Covid-19 situations, as the next question (Q23) highlights.

Q23 investigated the effect the Covid-19 pandemic had on the use of Hover. Above half of the respondents answered that they think they would have used Hover more if it had not been for the pandemic. The answers are essential since it confirms our perception of how the pandemic has influenced the results. The application should be tested when society is back to normal to get an improved depiction of how people’s habits are affected.

\textbf{Difference between the two groups} \\
The statements also revealed a considerable difference between the two groups. There was a statistically significant difference in four of the statements. The classmates were generally more positive. As explained in some of the previous subsections, the familiarity bias seems to be more apparent in the classmates’ answers. Their results are generally better than the student society.

The results also imply that the use of Hover among the group of classmates was more affected by the pandemic than the student society group. The results explain why the student society group used Hover more, while the group of classmates was more interested in continuing using a game like Hover.

\subsection{Summary}
In this chapter, we have discussed and compared the results from testing Hover on users. This included identifying synergies between the data collected on player engagement, motivation, enjoyment, and habits. This knowledge will be used in the next chapter to answer our research questions and conclude the project.
